\section{Обучение модели}

\subsection{Набор данных}

Обучение модели выполнено на синтетическом наборе данных. Набор данных сгенерирован под специфику решаемой задачи и включает `Train', `Val' и `Test'-коллекций csv-таблиц частотных характеристик и json-файлов с координатами полюсов/нулей \cite{gorbunov2026zerospolesmasks}.

Тренировочный набор данных представлен частотными характеристиками, соответствующими 72 случайным комбинациями полюсов/нулей. Каждой такой комбинации соответствуют 500 вариантов значений координат полюсов/нулей. Частотные характеристики вычислены для разных диапазонов частот по передаточной функции обобщенной модели линейного динамического звена:
\begin{equation*}
    H(s) = G \frac{\prod\limits_{m=1}^{M} \left(1 + \frac{s}{\omega_{\text{z}(m)}} \right)}{s^K \prod\limits_{n=1}^{N} \left(1 + \frac{s}{\omega_{\text{p}(n)}}\right)} e^{-s t_\text{delay}},
\end{equation*}
где
\begin{itemize}
    \item ${G}$ -- пропорциональный коэффициент, ${G \in \mathbb{R}}$;
    \item ${t_\text{delay}}$ -- задержка времени, ${t_\text{delay}>0}$;
    \item ${K}$ -- количество интеграторов, ${K \in \mathbb{N}}$;
    \item ${\omega_{\text{z}(m)}}$ -- частоты нулей передаточной функции ${H(s)}$ в количестве ${M}$, ${\omega_{\text{z}(m)} \in \mathbb{R}}$;
    \item ${\omega_{\text{p}(n)}}$ -- частоты полюсов передаточной функции ${H(s)}$ в количестве ${N}$, ${\omega_{\text{p}(n)} \in \mathbb{R}}$.
\end{itemize}

В тренировочном наборе данных параметры ${G}$, ${t_\text{delay}}$ заданы константами для последующего изменения в ходе обучения. Общий размер тренировочного набора данных -- 72x500=\num{36000} примеров.

В валидационном и тестовом наборах данных параметры ${G}$, ${t_\text{delay}}$ сгенерированы случайным образом. В тестовом наборе данных дополнительно частотные характеристики приближены к реальным данным путем введения шума с нормальным распределением. Размеры валидационного и тестового наборов -- \num{14400} и \num{7200} примеров, соответственно.

\subsection{Функция потерь}
Обучение производилось путем минимизации комбинированной функции потерь, вычисляемой по всем $N$ выходам декодера (глубокое супервизирование):
\begin{equation}\label{eqn:lossfunc}
    \mathcal{L} = \frac{1}{w} {\sum_{n=1}^{N} \Big[ w_n \Big( w^\text{BCE} \mathcal{L}_n^\text{BCE} + w^\text{Dice} \mathcal{L}_n^\text{Dice}} \Big) \Big],
\end{equation}
где:
\begin{itemize}
    \item ${\mathcal{L}^\text{BCE}}$ -- бинарная перекрестная энтропия;
    \item ${\mathcal{L}^\text{Dice}}$ -- функция потерь Дайса;
    \item ${w^\text{BCE}}$, ${w^\text{Dice}}$ -- весовые коэффициенты составляющих функции потерь.
\end{itemize}

Весовые коэффициенты ${w_n}$ выходов декодера заданы по возрастанию значений, так чтобы основному выходу модели соответствовало значение \num{1.0}. Весовые коэффициенты нормированы к сумме ${w = \sum_{n=1}^{N} w_n}$:
\begin{equation*}
    \displaystyle\sum_{n=1}^{N} \cfrac{w_n}{w} = 1.
\end{equation*}

Бинарная перекрестная энтропия:
\begin{equation*}
    \mathcal{L}^\text{BCE} = -\frac{1}{L} \sum_{i=1}^{L} \Big[ y_i \log\left(p_i \right) + \left(1 - y_i\right) \log \left( 1 - p_i \right) \Big],
\end{equation*}
где:
\begin{itemize}
    \item ${y_i}$ -- истинные (\enquote{ground truth}) значения маски, ${y_i \in \{0, 1\}}$;
    \item ${p_i}$ -- вероятности \eqref{eqn:logit2mask}, соответствующие логитам с выхода декодера, ${p_i \in \left[0; 1 \right]}$.
\end{itemize}

Функция потерь Дайса:
\begin{equation*}
    \mathcal{L}^{\text{Dice}} = 1 - 2 \frac{\sum_{i=1}^{L} d_i y_i}{\sum_{i=1}^{L} d_i + \sum_{i=1}^{L} y_i + \varepsilon}.
\end{equation*}

Бинарная перекрестная энтропия и функция потерь Дайса вычисляются по каждому каналу $\{1, C_\text{out}\}$; в функцию потерь \eqref{eqn:lossfunc} входят уже усредненные по каналам значения.

\subsection{Метрики}
В качестве основного показателя качества детекции модели использован коэффициент Дайса:
\begin{equation*}
    \text{Dice} = 2 \frac{\sum_{i=1}^{L} \hat{y}_i y_i}{\sum_{i=1}^{L} \hat{y}_i + \sum_{i=1}^{L} y_i + \varepsilon},
\end{equation*}
где:
\begin{itemize}
    \item ${y_i}$ -- истинные (\enquote{ground truth}) значения маски, ${y_i \in \{0, 1\}}$;
    \item ${\hat{y}_i}$ -- предсказания модели, ${\hat{y}_i \in \{0, 1\}}$.
\end{itemize}

К качестве вспомогательной метрики применен коэффициент Жаккара:
\begin{equation*}
    \text{IoU} = \frac{\sum_{i=1}^{L} \hat{y}_i y_i}{\sum_{i=1}^{L} \hat{y}_i + \sum_{i=1}^{L} y_i - \sum_{i=1}^{L} \hat{y}_i y_i + \varepsilon}.
\end{equation*}

Обучение выполнено на NVIDIA GeForce RTX3060 (12~ГБ), torch=2.6.0, CUDA=12.6, cuDNN=9.5.1 при следующих гиперпараметрах:
\begin{itemize}
    \item Весовые коэффициенты функции ошибки: ${w^\text{BCE}=\num{0.25}}$, ${w^\text{Dice}=\num{0.75}}$.
    \item Весовые коэффициенты глубокого супервизирования: ${w_1 = \num{1.00}}$, ${w_2 = \num{0.50}}$, ${w_3 = \num{0.25}}$.
    \item Размер батча: 360.
    \item Решатель: \enquote{Adam}.
\end{itemize}

Обучение модели выполнялось по лучшему коэффициенту Дайса на валидации ??????

\begin{table}[htbp]
    \caption{Метрики результирующие}
    \label{tbl:dice_iou}
    \begin{threeparttable}
        \begin{tabular}{|l|l|l|l|l|}
            \hline
                    &                & \textbf{Dice} & \textbf{IoU} & \textbf{Комментарий}                                  \\ \hline
            \textbf{1} & \textbf{Train} & 0.9092        & 0.8580       & Случайные аугментации: $G$, ${t_\text{delay}}$, шумы. \\ \hline
            \textbf{2} & \textbf{Val}   & 0.9393        & 0.9030       & Нет аугментаций.\tnote{*}                                  \\ \hline
            \textbf{3} & \textbf{Test}  & 0.9221        & 0.8774       & Нет аугментаций.\tnote{**}                                  \\ \hline
        \end{tabular}
        \begin{tablenotes}
            \small
            \item[*] Набор данных изначально сформирован с разными $G$, ${t_\text{delay}}$.
            \item[**] Набор данных изначально сформирован с шумами и разными $G$, ${t_\text{delay}}$.
        \end{tablenotes}
    \end{threeparttable}
\end{table}